

\documentclass[12pt]{report}
\usepackage[utf8]{inputenc}
\usepackage[russian]{babel}
\usepackage[14pt]{extsizes}
\usepackage{listings}
\usepackage{graphicx}
\usepackage{amsmath,amsfonts,amssymb,amsthm,mathtools} 
\usepackage{pgfplots}
\usepackage{filecontents}
\usepackage{float}
\usepackage{indentfirst}
\usepackage{eucal}
\usepackage{enumitem}
%s\documentclass[openany]{book}
\frenchspacing

\usepackage{titlesec}
\titleformat{\section}
{\normalsize\bfseries}
{\thesection}
{1em}{}
\titlespacing*{\chapter}{0pt}{-30pt}{8pt}
\titlespacing*{\section}{\parindent}{*4}{*4}
\titlespacing*{\subsection}{\parindent}{*4}{*4}

\usepackage{indentfirst} % Красная строка

\usetikzlibrary{datavisualization}
\usetikzlibrary{datavisualization.formats.functions}

\usepackage{amsmath}


% Для листинга кода:
\lstset{ %
	language=c,                 % выбор языка для подсветки (здесь это С)
	basicstyle=\small\sffamily, % размер и начертание шрифта для подсветки кода
	numbers=left,               % где поставить нумерацию строк (слева\справа)
	numberstyle=\tiny,           % размер шрифта для номеров строк
	stepnumber=1,                   % размер шага между двумя номерами строк
	numbersep=5pt,                % как далеко отстоят номера строк от подсвечиваемого кода
	showspaces=false,            % показывать или нет пробелы специальными отступами
	showstringspaces=false,      % показывать или нет пробелы в строках
	showtabs=false,             % показывать или нет табуляцию в строках
	frame=single,              % рисовать рамку вокруг кода
	tabsize=2,                 % размер табуляции по умолчанию равен 2 пробелам
	captionpos=t,              % позиция заголовка вверху [t] или внизу [b] 
	breaklines=true,           % автоматически переносить строки (да\нет)
	breakatwhitespace=false, % переносить строки только если есть пробел
	escapeinside={\#*}{*)}   % если нужно добавить комментарии в коде
}


\usepackage[left=2cm,right=2cm, top=2cm,bottom=2cm,bindingoffset=0cm]{geometry}
% Для измененных титулов глав:
\usepackage{titlesec, blindtext, color} % подключаем нужные пакеты
\definecolor{gray75}{gray}{0.75} % определяем цвет
\newcommand{\hsp}{\hspace{20pt}} % длина линии в 20pt
% titleformat определяет стиль
\titleformat{\chapter}[hang]{\Huge\bfseries}{\thechapter\hsp\textcolor{gray75}{|}\hsp}{0pt}{\Huge\bfseries}


% plot
\usepackage{pgfplots}
\usepackage{filecontents}
\usetikzlibrary{datavisualization}
\usetikzlibrary{datavisualization.formats.functions}

\begin{document}
	%\def\chaptername{} % убирает "Глава"
	\thispagestyle{empty}
	\begin{titlepage}
		\noindent \begin{minipage}{0.15\textwidth}
			\includegraphics[width=\linewidth]{b_logo}
		\end{minipage}
		\noindent\begin{minipage}{0.9\textwidth}\centering
			\textbf{Министерство науки и высшего образования Российской Федерации}\\
			\textbf{Федеральное государственное бюджетное образовательное учреждение высшего образования}\\
			\textbf{~~~«Московский государственный технический университет имени Н.Э.~Баумана}\\
			\textbf{(национальный исследовательский университет)»}\\
			\textbf{(МГТУ им. Н.Э.~Баумана)}
		\end{minipage}
		
		\noindent\rule{18cm}{3pt}
		\newline\newline
		\noindent ФАКУЛЬТЕТ $\underline{\text{«Информатика и системы управления»}}$ \newline\newline
		\noindent КАФЕДРА $\underline{\text{«Программное обеспечение ЭВМ и информационные технологии»}}$\newline\newline\newline\newline\newline
		
		\begin{center}
			\noindent\begin{minipage}{1.1\textwidth}\centering
				\Large\textbf{  Отчет по лабораторной работе №1}\newline
				\textbf{по дисциплине <<Функциональное и логическое}\newline
				\textbf{~~~программирование>>}\newline\newline
			\end{minipage}
		\end{center}
		
		\noindent\textbf{Тема} $\underline{\text{Списки в Lispе. Использование стандартных функций.}}$\newline\newline
		\noindent\textbf{Студент} $\underline{\text{Золотухин А. В.~~~~~~~~~~~~~~~~~~~~~~~~~~~~~~~~~~~~~~~~~~}}$\newline\newline
		\noindent\textbf{Группа} $\underline{\text{ИУ7-64Б~~~~~~~~~~~~~~~~~~~~~~~~~~~~~~~~~~~~~~~~~~~~~~~~~~}}$\newline\newline
		\noindent\textbf{Оценка (баллы)} $\underline{\text{~~~~~~~~~~~~~~~~~~~~~~~~~~~~~~~~~~~~~~~~~~~~~~~~~}}$\newline\newline
		\noindent\textbf{Преподаватели} $\underline{\text{Толпинская Н.Б., Строганов Ю. В.~~~~~~~~~~~~~~~~~~~~~~~~~~~~}}$\newline\newline\newline
		
		\begin{center}
			\vfill
			Москва~---~\the\year
			~г.
		\end{center}
	\end{titlepage}
	
	\chapter*{Теоретические вопросы}

	\section*{1.	Элементы языка: определение, синтаксис, представление в памяти.}
	
	Вся информация (данные и программы) в Lisp представляется в виде символьных выражений – S-выражений. 
	
	По определению: S-выражение ::= <атом> | <точечная пара>.
	
	Элементы языка Lisp включают в себя:
	\begin{itemize}
		\item Атомы:
		\begin{itemize} 
			\item символы (идентификаторы) – синтаксически – набор литер (букв и цифр), начинающихся с буквы;
			\item специальные символы – {T, Nil} (используются для обозначения логических констант);
			\item самоопределимые атомы – натуральные числа, дробные числа, вещественные числа, строки – последовательность символов, заключенных в двойные апострофы (например, “abc”);
		\end{itemize} 
		
		\item Более сложные данные – списки и точечные пары (структуры), которые строятся с помощью унифицированных структур – блоков памяти – бинарных узлов.
		Определения:
		
		Точечная пара ::= (<атом> . <атом>) | (<атом> . <точечная пара>) | (<точечная пара> . <атом>) | (<точечная пара> . <точечная пара>);
		
		Список ::= <пустой список> | <непустой список>, где 
		
		<пустой список> ::= () | Nil,
		
		<непустой список> ::= (<первый элемент> . <хвост>),
		
		<первый элемент> ::= <S-выражение>,
		
		<хвост> ::= <список>.
		
		Список – это частный случай S-выражения.
		
	\end{itemize}
	
	Синтаксически:
	
	любая структура (точечная пара или список) заключается в круглые скобки: (A . B) – точечная пара, (A) – список из одного элемента. Пустой список изображается как Nil или ();
	
	непустой список по определению может быть изображен: (A . (B . (C . (D . ())))),  допустимо изображение списка последовательностью атомов, разделенных пробелами – (A B C D).
	
	Элементы списка могут быть списками (любой список заключается в круглые скобки), например – (A (B C) (D С)). Таким образом, синтаксически наличие скобок является признаком структуры – списка или точечной пары.
	
	Любая непустая структура Lisp в памяти представляется списковой ячейкой, хранящей два указателя: на голову (первый элемент) и хвост – всё остальное. Точечная пара в памяти представляется бинарным узлом.
	
	\section*{2.	Особенности языка Lisp. Структура программы. Символ апостроф.}
	
	Отличительные особенности Lisp: только символьная обработка; все можно представить в виде функций.
	
	Вся информация (данные и программы) в Lisp представляется в виде символьных выражений – S-выражений. 
	
	По определению: S-выражение ::= <атом> | <точечная пара>.
	
	В зависимости от контекста одни и те же объекты могут играть роль переменных или констант, причем значения и того, и другого могут быть произвольной сложности. Если объект играет роль константы, то для объявления константы достаточно заблокировать его вычисление, то есть как бы взять его в кавычки, отмечающие буквально используемые фразы, не требующие обработки. 
	
	Апостроф – сокращённое обозначение функции quote.
	
	quote - блокирует вычисление своего аргумента. В качестве своего значения выдаёт сам аргумент, не вычисляя его. Перед константами - числами и атомами T, Nil можно не ставить апостроф.
	
	
	
	
	\section*{3.	Базис языка Lisp. Ядро языка.}
	Базис -- это минимальный набор инструментов языка и стркутур данных, который позволяет решить любые задачи.
	
	
	Базис Lisp :
	
	\begin{itemize}
		\item атомы и структуры (представляющиеся бинарными узлами);
		\item базовые (несколько) функций и функционалов: встроенные — примитивные 
		функции (atom, cons, list, car, cdr); специальные функции и функционалы (quote, lambda, defun).
	\end{itemize}
	
	Функцией называется правило, по которому каждому значению одного или нескольких  аргументов ставится в соответствие конкретное значение результата.
	
	Функционалом, или функцией высшего порядка называется функция, аргументом или  результатом которой является другая функция.
	
	Некоторые функции системы необходимо реализовывать в виде машинных подпрограмм, они образуют ядро системы. Ядро - основные действия, которые наиболее часто используются. Понятие <<ядро>> шире, чем понятие <<базис>>.
	
	
	\chapter*{Практические задания}	
	
	Практические задания на прикрепленном листе
	
	\bibliographystyle{utf8gost705u}  % стилевой файл для оформления по ГОСТу
	
	\bibliography{51-biblio}          % имя библиографической базы (bib-файла)
	
	
\end{document}